\chapter{Running, outputs and the water balance}\label{uch:outputs}
This chapter says how to run a coupled model, what Raven writes, and how to check that the water balance of a run is
right. Reference, Section~\ref{R-ch:outputs}, lists every output column, and Reference, Chapter~\ref{R-ch:c5}, the
accounts behind them.

\section{Running}
A coupled model runs like any Raven model:
\begin{lstlisting}
Raven Liard -o output/
\end{lstlisting}
Raven reads all input files, builds (or copies) the MODFLOW~6 model in \cmd{output/mf6/}, and then runs both models one
time step at a time. A successful run ends with \cmd{Exiting Gracefully: Successful Simulation}. Errors and warnings
are in \cmd{output/Raven\_errors.txt}.

Raven finds the MODFLOW~6 library by itself when it lies next to the Raven program or in \cmd{lib/mf6/} of the Raven
checkout. Otherwise name it with \cmd{:MF6Library path} in the \cmd{.rvg} file, or with the environment variable
\cmd{RAVEN\_MF6\_LIB}.

\section{The output files}
All rows are stamped with the date at the \emph{end} of the step, the moment the values refer to.
\begin{table}[!htbp]\centering
\caption{The groundwater outputs, beside Raven's usual files.}\label{tab:uoutputs}
\begin{tabularx}{\linewidth}{P{0.27\linewidth}Y}\toprule\hdr
File & What it holds\\\midrule
\cmd{GWModelSummary.txt} & the grid built (layers, cells, active cells, river cells), the MODFLOW library and its version, how each HRU is linked, run totals\\
\cmd{GWBudget.csv} & every step: recharge from Raven, seepage back to Raven, river exchange, storage change, MODFLOW's balance error, and every other exchange\\
\cmd{GWHRUState.csv} & every step: the water table and the depth to water under each coupled HRU\\
\cmd{GWHeads.csv} & the head at each monitoring well ($-1.2345$ when the screened cell is dry)\\
\cmd{GWConnectivity.csv} & the number of separate wet areas of the aquifer, and the wet share of each profile\\
\cmd{GWSoilExchange.csv} & with \cmd{:SoilZone PHYSICAL}: the soil zone of each HRU, every step\\
\cmd{GWHeads.nc} & with \cmd{:WriteNetCDFHeads}: the head in every cell (Raven built with NetCDF)\\
\cmd{Hydrographs.csv}, \cmd{Diagnostics.csv} & streamflow and fit statistics, as in any Raven run; \cmd{Diagnostics.csv} also scores the heads (\cmd{GW\_HEAD})\\
\cmd{WatershedStorage.csv} & Raven's balance; it includes all exchanged water\\
\cmd{solution.rvc} & the end state, with the groundwater blocks, for a restart (Chapter~\ref{uch:restarts})\\
\cmd{mf6/} & the MODFLOW~6 model Raven wrote or copied, its list files and the heads \cmd{gwf.hds}; it can be opened in ModelMuse or FloPy\\
\cmd{mf6/grid\_cells.geojson} & every cell as a longitude/latitude polygon, ready to be joined with heads in QGIS\\\bottomrule
\end{tabularx}
\end{table}

\section{Reading GWModelSummary.txt}
Read it after the first run of a new model. Check that:
\begin{itemize}
\item the grid and the number of active cells are what you expect;
\item every coupled HRU is linked, and its linked share is close to 1;
\item the polygon area of each HRU is close to its \cmd{.rvh} area;
\item with an existing model, the units, the packages taken over and kept, and the HRU coverage are right
(Chapter~\ref{uch:existing}).
\end{itemize}

\section{Reading GWBudget.csv}
\cmd{GWBudget.csv} has one row per step. Its columns are the aquifer's water account for that step, in m$^3$/d:
recharge in, seepage and river discharge out, river leakage in, and so on. \emph{Storage release} closes the account:
water released from the aquifer (negative: water stored).
\begin{table}[!htbp]\centering
\caption{The main columns of \cmd{GWBudget.csv}. Reference, Section~\ref{R-app:outputs}, lists them all.}\label{tab:ubudget}
\begin{tabularx}{\linewidth}{P{0.4\linewidth}Y}\toprule\hdr
Column & Meaning\\\midrule
time, date & model time [d] and date at the end of the step\\
Raven recharge, MF6 recharge & recharge leaving Raven, and entering MODFLOW\\
seepage to Raven, returned to Raven & drain discharge at the model top, and the volume added to Raven's stores\\
river leakage to aquifer, aquifer discharge to rivers & the river exchange, each direction\\
river loss carried to next step & a river loss the reach could not supply in the step (normally 0)\\
storage release & water released from aquifer storage (negative: stored)\\
MF6 balance error [m$^3$/d], [\%] & MODFLOW's balance error, and relative to the gross flow\\
outer iterations, converged & iterations used; 1 converged, 2 accepted at the relaxed tolerance, 0 not converged\\
water-table ET; wells specified, actual; GHB, CHD & the other exchanges (boundaries positive into the aquifer)\\
other model packages & with an existing model: net flow of the model's own packages\\
soil zone to aquifer, aquifer to soil zone, \dots & with \cmd{:SoilZone PHYSICAL}: the soil-zone exchange and its interface\\\bottomrule
\end{tabularx}
\end{table}

\begin{example}[: 15 June 1996 in the Liard run]
Recharge from Raven 21.4 million m$^3$/d and river leakage 3.0 million flowed in; seepage to Raven 19.1 million and
discharge to rivers 1.3 million flowed out. The aquifer stored the difference, 3.9 million m$^3$ (storage release
$-3\,925\,212$~m$^3$/d): a water-table rise of about 1~mm on average. Inflows minus outflows plus release is
0.00003~m$^3$/d, a relative error of about $10^{-10}$. Reference, Section~\ref{R-ch:balance}, shows Raven's side of the
same day.
\end{example}

\subsection{What a good run looks like}
\begin{enumerate}
\item \emph{MF6 balance error [\%]} stays well below 0.1\% every step. Larger values point to a step that did not
converge, or a head tolerance that is too loose (\cmd{:SolverHead\allowbreak Tolerance}, default 0.00001~m).
\item \emph{Raven recharge} and \emph{MF6 recharge} are equal, unless a recharge delay or fixed-head cells are used, or
recharge is rejected on saturated ground (\cmd{:SoilZone PHYSICAL}).
\item \emph{returned to Raven} equals \emph{seepage to Raven}.
\item \emph{converged} is 1 on every step (2: accepted at a relaxed tolerance; 0: not converged). Steps that did not
converge are flagged; see Chapter~\ref{uch:trouble}.
\item \emph{river loss carried to next step} is 0, except for a reach whose flow stopped within the step, a reservoir gain
MODFLOW did not supply (taken from the reach), or an existing model's stream package (Chapter~\ref{uch:existing}). The
carry is taken from the river the next step, and Raven warns.
\item With \cmd{:SoilZone PHYSICAL}: \emph{soil-zone shortfall} is zero and \emph{interface converged} is 1 on (nearly)
every step.
\item Raven's \cmd{WatershedStorage.csv} balance error stays at the first step's value (Raven's own, the same without
groundwater), apart from Raven's own known errors (Reference, Section~\ref{R-ch:limits}).
\end{enumerate}
Reference, Section~\ref{R-ch:balance}, gives the rules behind each item.

\section{The audit tool}\label{sec:uaudit}
\cmd{tools\_audit.py} checks all of these automatically. It lies in \cmd{examples/Liard\_groundwater/} and needs Python~3
with pandas:
\begin{lstlisting}
python tools_audit.py output
python ../Liard_groundwater/tools_audit.py output      # from example B's folder
\end{lstlisting}
Lines to look at:
\begin{itemize}
\item \cmd{recharge: (Raven - MF6 - dStore - toCHD)/Raven} should be 0;
\item \cmd{MF6 cumulative residual / gross flow} about $10^{-10}$ or smaller;
\item \cmd{unconverged steps} 0;
\item \cmd{ERROR lines} 0.
\end{itemize}
For example~B, a test run gave: recharge identity 0, \cmd{MF6 cumulative residual / gross flow} about $-1.9\times10^{-11}$,
\cmd{unconverged steps} 0, river bookkeeping about $-5\times10^{-12}$, \cmd{ERROR lines} 0 and \cmd{max outer iterations}
245 (the first step). With \cmd{:SoilZone PHYSICAL} the tool also closes every soil-zone identity.

\section{Heads everywhere, and particle tracking}
\begin{itemize}
\item \cmd{mf6/gwf.hds} (MODFLOW binary, readable with FloPy) holds the head in every cell every \cmd{:HeadSaveFrequency}
steps (default 30). \cmd{GWHeads.nc} holds them with \cmd{:WriteNetCDFHeads}.
\item With \cmd{:WriteGWBudgetFile} a generated model also saves MODFLOW's cell-by-cell flows (\cmd{mf6/gwf.cbc}) and the
porosity (\cmd{mf6/gwf.porosity}). \cmd{tools/make\_prt.py} then builds and runs a MODFLOW~6 particle-tracking model (PRT)
on the coupled flows, forward or backward from the wells (capture zones, travel times). PRT needs \cmd{:HeadSaveFrequency 1}.
Reference, Section~\ref{R-sec:budgetfile}, and Workshop~5 give the details.
\end{itemize}
